% =====================================================================
%  Chapter 21 — Numerical lessons
% =====================================================================
\chapter{What the Numerics Taught, Including the Hard Part}
\label{ch:numerical-lessons}

\begin{record}
\textbf{The hardest-won lesson in the book, stated once, without hedging.}

\smallskip
\emph{The \(c_K\) numerics were never discriminating.} Not
under-resolved, not inconclusive, not promising-but-limited: \textbf{incapable
of separating the hypothesis from its negation}, for a reason that was
available before the first run and that the programme itself proved in
\Sn{38}.

\smallskip
Four sessions --- \Sn{34}, \Sn{35}, \Sn{42}--\Sn{43}, \Sn{45} --- measured
\(c_K\) and its relatives. Every one of them returned a number. Not one of
them was evidence. The programme built the diagnostics carefully, ran them on
three initial-condition families at two resolutions and two viscosities,
swept the Reynolds number, corrected a resolution bias, corrected a misread
transient, calibrated the regression instrument on synthetic data, and probed
the one geometric corner the theory says matters --- and the whole of that
effort was, as evidence about the conjecture, worth exactly nothing.

\smallskip
Chapter~\ref{ch:layer4} gives the numbers. This chapter gives the reason, and
the three other lessons the numerical programme did produce.
\end{record}

\section{The non-discrimination argument}
\label{sec:nondiscrimination}

\subsection{The argument in four steps}

\begin{enumerate}[leftmargin=2.2em]
\item \texttt{conj:K-refined} is a \emph{conditional} statement. It asserts
  \(c_K\le C/\Lfac\) \emph{near a putative Type-I singularity}, in the limit
  \(\Lfac\to\infty\). It says nothing whatever about a flow that has no
  singularity.
\item A solver that remains globally regular is never in that regime. So
  neither the conjecture nor its negation predicts anything about the numbers
  such a solver produces. \emph{Both predict the same thing --- nothing.}
\item There is one reading on which a regular run could still have been a
  proxy: if the singular regime were merely a place the runs happened not to
  visit, a well-chosen run might approach it.
  \texttt{prop:regime-A-impossible} (line~6514) removes that reading. At a
  genuine blow-up time \(\limsup_{t\to T^-}\Omega(t)=\infty\); a resolved
  simulation is a bounded-enstrophy object throughout. The regime is not one
  the runs did not visit --- it is not of the \emph{type} a resolved run
  represents.
\item \Sn{38} proved, in the same proposition, that \(\sstar\)-decay is
  \emph{necessary} for blow-up given H1. A flow that never blows up therefore
  never exhibits the decay. It measures \(c_K\approx1\), and
  \(c_K\approx1\) is what \emph{both} hypotheses predict on any flow the
  solver can produce.
\end{enumerate}

A measurement that both sides of a disjunction predict cannot separate them.
That is the whole of Wall~6, and it depends on no number any run reported, on
no resolution, on no viscosity, and on no initial condition.

\begin{figure}[ht]
\centering
\begin{tikzpicture}[
  >={Latex[length=2mm]},
  bx/.style={draw, rounded corners=2pt, font=\scriptsize, align=center,
             inner sep=4pt, text width=36mm, minimum height=11mm},
  hy/.style={bx, draw=openblue, thick, fill=panelgrey},
  ms/.style={bx, draw=rulegrey},
  bad/.style={bx, draw=impossiblered, thick},
  ar/.style={->, thick}
]
\node[hy] (h) at (0,1.5)   {\textbf{H}: \(c_K\le C/\Lfac\)\\ near a Type-I
                            singularity};
\node[hy] (nh) at (0,-1.5) {\(\neg\)\textbf{H}: \(c_K\) does not decay\\
                            near a Type-I singularity};
\node[ms] (pred) at (5.6,0) {What each predicts for a\\
                             \emph{globally regular} solver};
\node[bad] (same) at (11.0,0)
  {\textbf{the same thing:}\\ nothing.\\
   Measured: \(c_K=0.92\)--\(2.24\)};

\draw[ar] (h) -- (pred);
\draw[ar] (nh) -- (pred);
\draw[ar] (pred) -- (same);

\node[bx, draw=provedgreen, very thick, text width=52mm] (s38) at (5.4,-3.2)
  {\Sn{38}, \texttt{prop:regime-A-impossible}: at any genuine blow-up
   \(\limsup\Omega=\infty\). A resolved run is bounded-enstrophy throughout,
   so it is not merely outside the regime --- it is not of that type};
\draw[ar, provedgreen, dashed] (s38.north) -- (pred.south);

\node[font=\scriptsize\itshape, text=impossiblered, align=center,
      text width=100mm] at (5.5,-4.8)
  {Four sessions of measurement. Zero bits of evidence. This was derivable
   before the first run.};
\end{tikzpicture}
\caption[Why the \(c_K\) measurements could not decide]{Wall~6 in diagram
form. The two hypotheses differ only on flows the solver cannot produce, and
agree on every flow it can. The measurements are real; their evidential value
is zero.}
\label{fig:nondiscrimination}
\end{figure}

\subsection{What each of the four sessions actually established}

The runs were not worthless as engineering. They were worthless as evidence
for the thing they were run for, and the two should be separated cleanly.

\begin{center}
\footnotesize
\begin{tabular}{@{}l p{0.34\textwidth} p{0.36\textwidth}@{}}
\toprule
Session & What it reported & What it established \\
\midrule
\Sn{34} & \(c_K\approx1\) in every flow; median \(1.06\)--\(1.94\); drifting
  toward \(1\) late-time; escalation lowered it
  & That the diagnostic runs and is stable. Nothing about the conjecture. \\[3pt]
\Sn{35} & \(c_K\) does not grow with Reynolds number over a fourfold viscosity
  range; slopes \(-0.031\)
  & That the diagnostic is not viscosity-sensitive in the accessible range.
    Nothing about the conjecture. \\[3pt]
\Sn{42}--\Sn{43} & Band concentration medians tame and resolution-robust;
  the spikes are an early transient, not concentration
  & A real correction to \Sn{42}'s own reading, and a real fact about
    transients. Nothing about the conjecture. \\[3pt]
\Sn{45} & \(\Lfac\)-slope of \(\log c_K\): \(-0.196\) (\(r^2=0.05\)) and
  \(+0.088\) (\(r^2=0.02\)) against the required \(-1\); (NB) fails outright
  in one run
  & That the decay's onset is not visible anywhere accessible. Which is what
    both hypotheses predict. \\
\bottomrule
\end{tabular}
\end{center}

\subsection{How the sessions themselves reported it}

The programme's own reporting was better than the programme's own strategy,
and that distinction is worth making because it is the difference between an
error of judgement and a failure of honesty.

\Sn{34} wrote that the measurement ``cannot yet distinguish \(c_K=O(1)\) from
the stronger conjectured \(c_K\le C/\Lfac\)''. \Sn{37} attached to a related
measurement the caveat that it ``measures ordinary transient vortex stretching
in a globally-regular, DECAYING flow \dots\ not the approach to an actual
singularity --- which by definition cannot be produced by a solver known to
stay regular''. \Sn{45}'s log records the negative slope as a negative signal
and adds, in the same paragraph, that ``the program has now been looking at
\(c_K\approx1\) for four sessions without ever seeing the decay the closure
needs''.

Every session named the limitation. No session drew the conclusion that the
limitation was fatal by construction. That inference waited for \Sn{46}, which
recorded it in one line --- ``\(c_K\) numerics were never discriminating'' ---
and for \Sn{47}, which turned it into Wall~6.

\begin{obsv}[The specific shape of the mistake]
\label{obs:conditional-proxy}
The error was not measuring the wrong quantity. \(c_K\) is exactly the
quantity the conjecture is about, computed exactly as the conjecture defines
it. The error was measuring \emph{a conditional statement's consequent on flows
that do not satisfy its antecedent}, and treating agreement as support.

\smallskip
The general form: before running a diagnostic for a conditional claim, write
down what the claim's \emph{negation} predicts for the flows you can actually
produce. If the two predictions coincide, the diagnostic is not an experiment.
This costs one paragraph and would have saved four sessions.
\end{obsv}

\begin{obsv}[What Wall 6 does not indict]
\label{obs:wall6-scope}
It indicts a family of measurements, not the layer, and the document is
careful about the line. \emph{Permitted:} measurements whose answer differs
between competing mechanisms \emph{in a regular flow} --- questions about a
mechanism, not proxies for blow-up. \emph{Forbidden:} computing \(c_K\),
\(c_K^{\mathrm{band}}\) or the band ratios again \emph{as evidence} for or
against the conjectures. They may be computed as context.

\smallskip
Several Layer-4 results are of the permitted kind and are genuine: the
\(\varepsilon\)-independence of \(\delta\), which is a fact about the
equations; the sampling-bias discovery; the resolution correction to \(c_K\);
the early-transient explanation of the band spikes; and the demonstration that
free viscous decay alone does not suppress gradient concentration.
\end{obsv}

\section{Resolution: three ways it lied}
\label{sec:resolution}

\subsection{Systematic suppression, found by refining}

The \(N=64\) Taylor--Green \(c_K\) values were resolution-suppressed by
roughly a factor of two. At \(N=64\) the spectral tail fraction ranged over
\(0.07\)--\(0.59\); at \(N=128\) it fell to
\(1.3\times10^{-5}\)--\(2.6\times10^{-5}\), and the measured \(c_K\) rose from
\(\approx1.05\) to \(\approx2.24\). That is systematic spectral pile-up at the
grid cutoff, not noise, and it moved Taylor--Green from ``beside the
adversarial series'' to ``noticeably above it''.

A further probe, \texttt{EXP-L4-KCORR-NU-TG-104}, showed the tail fraction
climbing back to \(0.0082\) late in the Taylor--Green cascade \emph{even at}
\(N=128\).

\begin{record}
\textbf{The rule this produced, and it was never discharged.} Resolution
adequacy is \emph{time-window-dependent}, not a property of \((N,\nu)\) alone.
A run can be well resolved for the first half of its trajectory and marginal
for the second. Reporting a single tail fraction per run therefore conceals
exactly the regime --- late-time, high-\(M\) --- that every one of the
programme's conjectures is about. \Sn{36} flagged this and it was never
resolved; an \(N=256\) check was named as the test if the exact magnitude of
\(c_K\) ever became analytically load-bearing. It did not, because of Wall~6.
\end{record}

\subsection{A hard grid floor masquerading as a physical threshold}

\Sn{40}'s anisotropic-collapse sweep produced the programme's only \textsf{FAIL}
verdict in Layer~3: \texttt{EXP-L3-R2-ANISO0050-JAX-128}, at aspect ratio
\(0.05\), with \(\delta_{\max}=0.00\) exactly. The run was not truncated ---
\(\theta\) stayed non-negative, energy decayed normally from
\(1.328\times10^{-3}\) to \(7.63\times10^{-4}\), and it completed all \(69\)
steps.

Before treating that as a finding, the session checked for an artefact. At
\(N=128\), the module's axial length \(\ell_z\) hits a hard grid floor: the
requested \(0.05\times\ell_r=0.0393\) is below one grid cell, so it is clamped
to exactly \(2\pi/128=0.0491\). The intended core geometry is representable as
literally one grid point.

A control isolated the explanation.

\begin{center}
\small
\begin{tabular}{@{}l r r l r l@{}}
\toprule
\texttt{exp\_id} & \(N\) & aspect & floor-clamped? & \(\delta_{\max}\) & verdict \\
\midrule
\texttt{EXP-L3-R2-ANISO0100-JAX-064} & \(64\) & \(0.10\)
  & yes (\(\ell_z=2\pi/64\)) & \(0.50\) & \textsf{PASS} \\
\texttt{EXP-L3-R2-ANISO0100-JAX-128} & \(128\) & \(0.10\)
  & no (\(\ell_z=0.0785\)) & \(0.50\) & \textsf{PASS} \\
\texttt{EXP-L3-R2-ANISO0050-JAX-128} & \(128\) & \(0.05\)
  & \textbf{yes} (\(\ell_z=2\pi/128\)) & \(\mathbf{0.00}\) & \textsf{FAIL} \\
\bottomrule
\end{tabular}
\end{center}

\medskip
Aspect ratio \(0.10\) is genuinely resolution-converged --- identical
\(\delta_{\max}\) at two resolutions. The sole \textsf{FAIL} coincides exactly
with the one configuration pinned to the floor. The session's conclusion is
stated with the right strength and no more: ``most likely a discretization
artefact \dots\ This is inference from a single controlled comparison, not
proof the floor explains everything.'' Cleanly testing the real \(0.05\)
geometry needs \(N\gtrsim512\), which was not attempted.

The honest state of that sweep is therefore: \(\delta_{\max}>0\) at every
\emph{properly resolved} aspect ratio tested, down to \(0.1\), confirmed at two
resolutions; and no genuine adversarial anisotropy threshold found --- only a
hard numerical resolution floor.

\begin{obsv}[Design the diagnostic so its own limits are visible]
\label{obs:floor-visible}
The floor-clamp was catchable only because the module recorded the
\emph{requested} and \emph{realised} \(\ell_z\) separately, so a human could
see that they differed. Had it silently clamped, the \textsf{FAIL} would have
entered the record as an anisotropy threshold --- a physical finding, and a
false one, in the single most adversarial experiment the programme ran.
\end{obsv}

\subsection{An early transient mistaken for a resolution artefact, and then
correctly identified}

\Sn{42} ran the band-concentration diagnostic at \(N=64\) and found the
medians tame but two maxima extreme: \(\mathrm{conc}_D=63.1\) for
Taylor--Green and \(192.8\) for shear. It flagged them as ``likely resolution
artefacts'', explicitly citing the programme's own documented \(N=64\)
under-resolution history --- flagged, not explained away.

\Sn{43} refined to \(N=128\), at about \(2.7\) hours per run, and found the
reading \emph{half wrong}. Taylor--Green's spikes \emph{grew} fivefold under
refinement; shear's shrank. A statistic that moves in opposite directions under
refinement on two initial conditions is not converged, and nothing can rest on
it.

The decisive analysis was neither more resolution nor more runs: it was
splitting each recorded timeseries at \(t=3\). Across \emph{all} runs, both
resolutions, all three initial-condition families, \emph{every} maximum occurs
at \(t<3\). Taylor--Green at \(N=128\) --- the largest spike in the whole
study, \(135.95\) --- never exceeds \(0.63\) anywhere in the bulk, and its bulk
maxima agree with the \(N=64\) run to within \(7\)--\(9\%\).

The accurate characterisation is an \emph{early-transient} effect: the initial
conditions are still smooth and symmetric, the high-vorticity set has not been
developed by nonlinear mixing, and band statistics over few unrepresentative
points are volatile. That is the wrong time regime to bear on a conjecture
about the approach to a singularity. \Sn{43} also notes what it did
\emph{not} establish: which of two plausible mechanisms produces the
volatility. Both are offered as hypotheses, not claims.

\begin{obsv}[Two corrections, and the second one corrects the first]
\label{obs:s43-self-correction}
\Sn{43}'s own Result~3 said bulk maxima ``converge downward under
refinement''. Its addendum corrects that: the generalisation came from shear
alone, and Taylor--Green's bulk maxima in fact \emph{rose} slightly,
\(1.57\to1.71\). The accurate statement is that bulk statistics are
resolution-\emph{stable}, moving by under \(10\%\), and where they moved
substantially they moved downward.

\smallskip
Nothing in the conclusion depended on the direction --- only on the bulk being
small and stable. But the correction was made anyway, in the same document, in
the same session. That is the standard this book is trying to record.
\end{obsv}

\section{Sampling: the forty-seven orders of magnitude}
\label{sec:sampling}

Chapter~\ref{ch:layer4} tells this story with the numbers; it belongs here as
a lesson because the failure mode is not specific to fluid dynamics.

The mean-Morrey quantity is a \emph{supremum} over balls. The first
implementation sampled eight ball centres on a fixed grid. For Taylor--Green
and shear fields --- built from sines and cosines --- those centres land on
symmetry nodes where the mean vanishes identically by construction. The
measured maximum was \(7.70\times10^{-47}\). With \(64\) random centres and a
logged seed, it was \(11.03\).

\begin{record}
\textbf{The rule.} When the quantity is a supremum over a family, sample the
family randomly with a logged seed --- never on a lattice that shares the
field's symmetries. The programme's later diagnostics sample at the vorticity
maximum rather than on a grid, for the same reason.

\smallskip
\textbf{Why the comparison was possible at all:} both sets of rows were
retained in the database. A programme that had overwritten the bad rows with
the good ones would have the right numbers and no record that the failure
mode exists.
\end{record}

\section{Calibrate the instrument on planted data}
\label{sec:calibration}

\Sn{45}'s regression of \(\log c_K\) on \(\log\Lfac\) is the one place the
programme's answer depended on the behaviour of a statistical instrument, and
it is the one place the instrument was calibrated before use.

\texttt{l\_dependence\_summary} was tested against \emph{planted synthetic
power laws} and recovers slopes of \(-1.0\) and \(-0.5\) exactly, with
\(r^2=1\). \texttt{sobolev\_H3\_norm} was checked against its analytic value
on a single Fourier mode.

This matters for how the negative result is read. When the instrument then
returns \(-0.196\) with \(r^2=0.05\), the low \(r^2\) cannot be dismissed as a
defect of the fit --- the fit is known to be exact on data of the shape being
looked for. The low \(r^2\) \emph{is} the finding: across
\(\Lfac\in[1.35,4.68]\), a \(3.5\times\) span over which a clean \(C/\Lfac\)
law would force a \(3.5\times\) drop, \(c_K\) stays pinned between \(0.92\) and
\(1.21\).

The same session records two design decisions of the same character, both
worth generalising:

\begin{itemize}[leftmargin=1.6em]
\item \emph{Hold energy fixed and sweep a ratio.} Sweeping raw strain
  amplitude would have confounded ``more strain'' with ``more energy'' and
  moved every diagnostic for a trivial reason. The swept parameter is the
  strain-to-swirl ratio at fixed energy.
\item \emph{The tilt perturbation is mandatory, not cosmetic.} A perfectly
  aligned axisymmetric jet has \(\ehat\) constant, so \(w=\abs{\nabla\ehat}^2\equiv0\)
  and every correlation diagnostic collapses to its zero-by-convention branch.
  An experiment that returns zeros because its target quantity is identically
  zero is not a null result.
\end{itemize}

\begin{obsv}[Calibration is what makes a negative result usable]
\label{obs:calibration-value}
The \(\Lfac\)-decay measurement is, by Wall~6, not evidence about the
conjecture. It is nonetheless the best-executed measurement in the programme,
and the reason is the calibration: it is the only result whose \emph{null}
reading can be trusted, because the instrument was shown in advance to detect
the signal it failed to find.

\smallskip
That is the general point. A positive result survives a poorly characterised
instrument; a negative one does not. The programme learned this at the end,
and applied it exactly once.
\end{obsv}

\section{A diagnostic whose regime does not match its conjecture's}
\label{sec:gamma}

\texttt{EXP-L4-GAM-ADV-001} is the programme's most-quoted negative row and
its \texttt{verdict} field is a sentence rather than a verdict:
``\textsf{NOT SUPPRESSED: NS evolution does not increase Gamma exponent
(0.000 -> -3.263). Possible gap in Claim A conjecture for adversarial ICs.}''

The reading in the \Sn{27} log is the correct one. The evolution in that run
is \emph{free viscous decay}, \(\partial_tu=\nu\Delta u\), with the nonlinear
term omitted --- and the mechanism the conjecture appeals to lives entirely in
the nonlinearity that was omitted. The run does not falsify the conjecture; it
establishes that linear diffusion alone does not suppress gradient
concentration, which nobody had claimed. The stated next step, rerunning the
adversarial data through the full spectral solver, was never performed under
that identifier, and the row stands in the database unresolved.

Chapter~\ref{ch:inheritance} calls this a prefiguring of Wall~6, and it is the
same disease in a milder form: \emph{a diagnostic whose regime does not match
the conjecture's cannot test it}. In \Sn{27} the mismatch was an omitted term
and could have been fixed by rerunning. In \Sn{34}--\Sn{45} the mismatch was
between a regular flow and a singular one, and could not be fixed at all.

\section{The lessons, collected}
\label{sec:numerics-lessons-table}

\begin{table}[ht]
\centering
\footnotesize
\begin{tabular}{@{}p{0.30\textwidth} p{0.28\textwidth} p{0.34\textwidth}@{}}
\toprule
Lesson & Where it was learned & Cost of not knowing it \\
\midrule
Before running a diagnostic for a conditional claim, write down what the
claim's \emph{negation} predicts on the flows you can produce
  & \Sn{46}, \Sn{47}; provable from \Sn{38}
  & Four sessions of measurement with zero evidential value \\[3pt]
Resolution adequacy is time-window-dependent
  & \Sn{36}, \Sn{37}
  & A factor of two, systematic, in the headline statistic \\[3pt]
Record requested \emph{and} realised parameters so a clamp is visible
  & \Sn{40}
  & A grid floor would have entered the record as a physical threshold \\[3pt]
Split a timeseries by regime before believing an extremum
  & \Sn{43}
  & A \(214\times\) spike read as concentration rather than transient \\[3pt]
Sample a supremum randomly, with a logged seed
  & \Sn{25}
  & \(47\) orders of magnitude of bias \\[3pt]
Calibrate a fitting instrument on planted data before trusting a null
  & \Sn{45}
  & A null result that cannot be distinguished from a broken fit \\[3pt]
Sweep a ratio at fixed energy, not an amplitude
  & \Sn{45}
  & Every diagnostic moving for a trivial reason \\[3pt]
Retain superseded rows in the database
  & \Sn{25}, \Sn{36}
  & No record that the failure mode exists \\[3pt]
Match the diagnostic's regime to the conjecture's
  & \Sn{27}, and again \Sn{34}--\Sn{45}
  & An unresolved sentence in a verdict field; then Wall~6 \\
\bottomrule
\end{tabular}
\caption[The numerical lessons]{Nine lessons, and what each cost. Only the
first is fatal: the others are corrigible errors that the programme did in
fact correct.}
\label{tab:numerical-lessons}
\end{table}

\begin{obsv}[The honest verdict on the numerical programme]
\label{obs:numerics-verdict}
Roughly \(124\) logged experiments across four layers, thirty-one database
tables, and something over a thousand tests. The Layer-1 to Layer-3 work is
sound and did what it was built for: it validated the property engine,
confirmed \(A(T)>0\) across every fluid and every profile, verified the
\(\mu(T)\to\nu\) limit in two and three dimensions, and searched for blow-up
at three resolutions without finding it. Layer~4 produced five results of
genuine mechanism-level interest.

\smallskip
And the single family of measurements the later analytical programme most
needed --- the one that ran for four sessions and consumed the largest share
of the numerical effort after \Sn{33} --- was incapable of answering the
question it was built for, on grounds the programme itself had proved. That is
the numerical programme's headline, and it is not a resolution problem, a
compute problem, or a modelling problem. It is a problem of not having asked,
in advance, what the negation predicted.
\end{obsv}
